\subsection{device}\label{device}
Das \verb|device| Modul stellt eine generische Geräte\-Abstraktion für
Byte\-orientierte Ein\-/Ausgabe bereit. Ein Gerät wird durch \verb|device_t|
(Listing \ref{device_t}) beschrieben -- eine Tabelle aus Funktionszeigern
für \verb|open|, \verb|read|, \verb|write|, \verb|ioctrl|, \verb|close| und
einen \verb|ready_cb|\-Rückruf, ergänzt um anwendungsspezifische
\verb|user_data|. Konkrete Geräte registrieren sich mit \verb|device_init()|
und werden anschließend über ein \verb|dev_handle_t| ($\ge 0$ gültig)
angesprochen; die \verb|device_*|\-Funktionen prüfen das Handle und delegieren
an die hinterlegten Zeiger. Genutzt wird das Modul u.\,a.\ vom \verb|serial|
(Abschnitt \ref{serial}, \verb|serial_io|).

\begin{listing}
\begin{minted}{C}
typedef struct device_s {
    void *user_data;
    em_msg (*open)(dev_handle_t, void *user_data);
    em_msg (*read)(dev_handle_t, uint8_t *buffer, int16_t *cnt);
    em_msg (*write)(dev_handle_t, const uint8_t *buffer, int16_t cnt);
    em_msg (*ioctrl)(dev_handle_t, dev_command_t cmd, int16_t value);
    em_msg (*close)(dev_handle_t);
    em_msg (*ready_cb)(device_state_t state);
    uint8_t dev_type;
} device_t;
\end{minted}
\caption{Geräte\-Schnittstelle \texttt{device\_t}}
\label{device_t}
\end{listing}

\noindent Welche Funktionen ein Gerät tatsächlich anbietet, wird über das
Bitfeld \verb|dev_func_t| (\verb|DEV_OPEN|, \verb|DEV_READ|, \verb|DEV_WRITE|,
\dots) ausgedrückt. \verb|device_check()| prüft, ob die durch eine solche
Maske geforderten Funktionszeiger gesetzt sind.

\subsubsection*{API}
{\sloppy\emergencystretch=3em
\begin{description}
\item[\texttt{device\_init(dev, user\_data)}] Registriert das Gerät und ruft
  dessen \verb|open| auf. Rückgabe: \verb|dev_handle_t| (zugeteiltes Handle,
  $0$ bei Fehler bzw.\ voller Tabelle).
\item[\texttt{device\_check(hdl, dev\_type)}] Prüft, ob die per Maske
  \verb|dev_type| geforderten Funktionen vorhanden sind. Rückgabe:
  \verb|em_msg| (\verb|EM_OK|, wenn alle geforderten Zeiger gesetzt sind,
  sonst \verb|EM_ERR|).
\item[\texttt{device\_read/write(hdl, buffer, cnt)}] Lesen bzw.\ schreiben über
  die Funktionszeiger des Geräts. Rückgabe: \verb|em_msg| (\verb|EM_ERR|, wenn
  das Gerät oder der jeweilige Zeiger fehlt).
\item[\texttt{device\_reset(hdl)}] Setzt das Gerät auf \verb|DEVICE_RESET|
  (alle Zeiger \verb|NULL|) zurück. Rückgabe: \verb|em_msg| (\verb|EM_OK|,
  bzw.\ \verb|EM_ERR| bei unbekanntem Gerät).
\item[\texttt{device\_print(hdl)}] Gibt die hinterlegten Zeiger zu
  Debug\-Zwecken aus. Rückgabe: \verb|void|.
\end{description}
\par}
